\documentclass[12pt, a4paper]{article}

\usepackage[a4paper, margin=1in]{geometry}

\usepackage[utf8]{inputenc}
\usepackage[mathscr]{euscript}
\usepackage{amssymb,amsmath,amsthm,amsfonts}
\usepackage[shortlabels]{enumitem}
\usepackage{multicol,multirow}
\usepackage{lipsum}
\usepackage{balance}
\usepackage{calc}
\usepackage[colorlinks=true,citecolor=blue,linkcolor=blue]{hyperref}
\usepackage{import}
\usepackage{xifthen}
\usepackage{pdfpages}
\usepackage{transparent}
\usepackage{listings}

\newcommand{\incfig}[2][1.0]{
    \def\svgwidth{#1\columnwidth}
    \import{./figures/}{#2.pdf_tex}
}

\newcommand{\incimg}[2][1.0]{
  \includegraphics[width=#1\columnwidth]{./figures/#2}
}

\newlist{enumproof}{enumerate}{4}
\setlist[enumproof,1]{label=\arabic*., parsep=1em}
\setlist[enumproof,2]{label=\arabic{enumproofi}.\arabic*., parsep=1em}
\setlist[enumproof,3]{label=\arabic{enumproofi}.\arabic{enumproofii}.\arabic*., parsep=1em}
\setlist[enumproof,4]{label=\arabic{enumproofi}.\arabic{enumproofii}.\arabic{enumproofiii}.\arabic*., parsep=1em}

\renewcommand{\qedsymbol}{\ensuremath{\blacksquare}}

\lstdefinestyle{mystyle}{
  language=Python, % Set the language to Python
  commentstyle=\color{codegray}\itshape, % Python comments (#)
  keywordstyle=\color{orange}\bfseries,  % Python reserved keywords
  stringstyle=\color{mauve},             % Strings
  basicstyle={\ttfamily\footnotesize},   % Monospace font
  breakatwhitespace=true,         
  breaklines=true,                 
  captionpos=b,                    
  keepspaces=true,                 
  numbers=none,                    
  tabsize=2,
  showstringspaces=false,
  % Extra keyword groups for Python
  morekeywords=[2]{self, cls, __init__, __str__, __repr__}, 
  keywordstyle=[2]\color{identblue},     % Special identifiers / methods
  morekeywords=[3]{int, float, bool, str, list, dict, set, tuple, object, None, True, False},
  keywordstyle=[3]\color{codegreen},     % Built-in types
  morekeywords=[4]{print, input, open, range, len, type, isinstance, super, enumerate, zip, map, filter, sum, min, max, sorted},
  keywordstyle=[4]\color{cyan},          % Built-in functions
}

% Things Lie
\newcommand{\kb}{\mathfrak b}
\newcommand{\kg}{\mathfrak g}
\newcommand{\kh}{\mathfrak h}
\newcommand{\kn}{\mathfrak n}
\newcommand{\ku}{\mathfrak u}
\newcommand{\kz}{\mathfrak z}
\DeclareMathOperator{\Ext}{Ext} % Ext functor
\DeclareMathOperator{\Tor}{Tor} % Tor functor
\newcommand{\gl}{\opname{\mathfrak{gl}}} % frak gl group
\renewcommand{\sl}{\opname{\mathfrak{sl}}} % frak sl group chktex 6

% More script letters etc.
\newcommand{\SA}{\mathcal A}
\newcommand{\SB}{\mathcal B}
\newcommand{\SC}{\mathcal C}
\newcommand{\SF}{\mathcal F}
\newcommand{\SG}{\mathcal G}
\newcommand{\SH}{\mathcal H}
\newcommand{\OO}{\mathcal O}

\newcommand{\SCA}{\mathscr A}
\newcommand{\SCB}{\mathscr B}
\newcommand{\SCC}{\mathscr C}
\newcommand{\SCD}{\mathscr D}
\newcommand{\SCE}{\mathscr E}
\newcommand{\SCF}{\mathscr F}
\newcommand{\SCG}{\mathscr G}
\newcommand{\SCH}{\mathscr H}

% Mathfrak primes
\newcommand{\km}{\mathfrak m}
\newcommand{\kp}{\mathfrak p}
\newcommand{\kq}{\mathfrak q}

% number sets
\newcommand{\RR}[1][]{\ensuremath{\ifstrempty{#1}{\mathbb{R}}{\mathbb{R}^{#1}}}}
\newcommand{\NN}[1][]{\ensuremath{\ifstrempty{#1}{\mathbb{N}}{\mathbb{N}^{#1}}}}
\newcommand{\ZZ}[1][]{\ensuremath{\ifstrempty{#1}{\mathbb{Z}}{\mathbb{Z}^{#1}}}}
\newcommand{\QQ}[1][]{\ensuremath{\ifstrempty{#1}{\mathbb{Q}}{\mathbb{Q}^{#1}}}}
\newcommand{\CC}[1][]{\ensuremath{\ifstrempty{#1}{\mathbb{C}}{\mathbb{C}^{#1}}}}
\newcommand{\PP}[1][]{\ensuremath{\ifstrempty{#1}{\mathbb{P}}{\mathbb{P}^{#1}}}}
\newcommand{\HH}[1][]{\ensuremath{\ifstrempty{#1}{\mathbb{H}}{\mathbb{H}^{#1}}}}
\newcommand{\FF}[1][]{\ensuremath{\ifstrempty{#1}{\mathbb{F}}{\mathbb{F}^{#1}}}}
% expected value
\newcommand{\EE}{\ensuremath{\mathbb{E}}}
\newcommand{\charin}{\text{ char }}
\DeclareMathOperator{\sign}{sign}
\DeclareMathOperator{\Aut}{Aut}
\DeclareMathOperator{\Inn}{Inn}
\DeclareMathOperator{\Syl}{Syl}
\DeclareMathOperator{\Gal}{Gal}
\DeclareMathOperator{\GL}{GL} % General linear group
\DeclareMathOperator{\SL}{SL} % Special linear group

%---------------------------------------
% BlackBoard Math Fonts :-
%---------------------------------------

%Captital Letters
\newcommand{\bbA}{\mathbb{A}}	\newcommand{\bbB}{\mathbb{B}}
\newcommand{\bbC}{\mathbb{C}}	\newcommand{\bbD}{\mathbb{D}}
\newcommand{\bbE}{\mathbb{E}}	\newcommand{\bbF}{\mathbb{F}}
\newcommand{\bbG}{\mathbb{G}}	\newcommand{\bbH}{\mathbb{H}}
\newcommand{\bbI}{\mathbb{I}}	\newcommand{\bbJ}{\mathbb{J}}
\newcommand{\bbK}{\mathbb{K}}	\newcommand{\bbL}{\mathbb{L}}
\newcommand{\bbM}{\mathbb{M}}	\newcommand{\bbN}{\mathbb{N}}
\newcommand{\bbO}{\mathbb{O}}	\newcommand{\bbP}{\mathbb{P}}
\newcommand{\bbQ}{\mathbb{Q}}	\newcommand{\bbR}{\mathbb{R}}
\newcommand{\bbS}{\mathbb{S}}	\newcommand{\bbT}{\mathbb{T}}
\newcommand{\bbU}{\mathbb{U}}	\newcommand{\bbV}{\mathbb{V}}
\newcommand{\bbW}{\mathbb{W}}	\newcommand{\bbX}{\mathbb{X}}
\newcommand{\bbY}{\mathbb{Y}}	\newcommand{\bbZ}{\mathbb{Z}}

%---------------------------------------
% MathCal Fonts :-
%---------------------------------------

%Captital Letters
\newcommand{\mcA}{\mathcal{A}}	\newcommand{\mcB}{\mathcal{B}}
\newcommand{\mcC}{\mathcal{C}}	\newcommand{\mcD}{\mathcal{D}}
\newcommand{\mcE}{\mathcal{E}}	\newcommand{\mcF}{\mathcal{F}}
\newcommand{\mcG}{\mathcal{G}}	\newcommand{\mcH}{\mathcal{H}}
\newcommand{\mcI}{\mathcal{I}}	\newcommand{\mcJ}{\mathcal{J}}
\newcommand{\mcK}{\mathcal{K}}	\newcommand{\mcL}{\mathcal{L}}
\newcommand{\mcM}{\mathcal{M}}	\newcommand{\mcN}{\mathcal{N}}
\newcommand{\mcO}{\mathcal{O}}	\newcommand{\mcP}{\mathcal{P}}
\newcommand{\mcQ}{\mathcal{Q}}	\newcommand{\mcR}{\mathcal{R}}
\newcommand{\mcS}{\mathcal{S}}	\newcommand{\mcT}{\mathcal{T}}
\newcommand{\mcU}{\mathcal{U}}	\newcommand{\mcV}{\mathcal{V}}
\newcommand{\mcW}{\mathcal{W}}	\newcommand{\mcX}{\mathcal{X}}
\newcommand{\mcY}{\mathcal{Y}}	\newcommand{\mcZ}{\mathcal{Z}}


%---------------------------------------
% Bold Math Fonts :-
%---------------------------------------

%Captital Letters
\newcommand{\bmA}{\boldsymbol{A}}	\newcommand{\bmB}{\boldsymbol{B}}
\newcommand{\bmC}{\boldsymbol{C}}	\newcommand{\bmD}{\boldsymbol{D}}
\newcommand{\bmE}{\boldsymbol{E}}	\newcommand{\bmF}{\boldsymbol{F}}
\newcommand{\bmG}{\boldsymbol{G}}	\newcommand{\bmH}{\boldsymbol{H}}
\newcommand{\bmI}{\boldsymbol{I}}	\newcommand{\bmJ}{\boldsymbol{J}}
\newcommand{\bmK}{\boldsymbol{K}}	\newcommand{\bmL}{\boldsymbol{L}}
\newcommand{\bmM}{\boldsymbol{M}}	\newcommand{\bmN}{\boldsymbol{N}}
\newcommand{\bmO}{\boldsymbol{O}}	\newcommand{\bmP}{\boldsymbol{P}}
\newcommand{\bmQ}{\boldsymbol{Q}}	\newcommand{\bmR}{\boldsymbol{R}}
\newcommand{\bmS}{\boldsymbol{S}}	\newcommand{\bmT}{\boldsymbol{T}}
\newcommand{\bmU}{\boldsymbol{U}}	\newcommand{\bmV}{\boldsymbol{V}}
\newcommand{\bmW}{\boldsymbol{W}}	\newcommand{\bmX}{\boldsymbol{X}}
\newcommand{\bmY}{\boldsymbol{Y}}	\newcommand{\bmZ}{\boldsymbol{Z}}
%Small Letters
\newcommand{\bma}{\boldsymbol{a}}	\newcommand{\bmb}{\boldsymbol{b}}
\newcommand{\bmc}{\boldsymbol{c}}	\newcommand{\bmd}{\boldsymbol{d}}
\newcommand{\bme}{\boldsymbol{e}}	\newcommand{\bmf}{\boldsymbol{f}}
\newcommand{\bmg}{\boldsymbol{g}}	\newcommand{\bmh}{\boldsymbol{h}}
\newcommand{\bmi}{\boldsymbol{i}}	\newcommand{\bmj}{\boldsymbol{j}}
\newcommand{\bmk}{\boldsymbol{k}}	\newcommand{\bml}{\boldsymbol{l}}
\newcommand{\bmm}{\boldsymbol{m}}	\newcommand{\bmn}{\boldsymbol{n}}
\newcommand{\bmo}{\boldsymbol{o}}	\newcommand{\bmp}{\boldsymbol{p}}
\newcommand{\bmq}{\boldsymbol{q}}	\newcommand{\bmr}{\boldsymbol{r}}
\newcommand{\bms}{\boldsymbol{s}}	\newcommand{\bmt}{\boldsymbol{t}}
\newcommand{\bmu}{\boldsymbol{u}}	\newcommand{\bmv}{\boldsymbol{v}}
\newcommand{\bmw}{\boldsymbol{w}}	\newcommand{\bmx}{\boldsymbol{x}}
\newcommand{\bmy}{\boldsymbol{y}}	\newcommand{\bmz}{\boldsymbol{z}}

%---------------------------------------
% Scr Math Fonts :-
%---------------------------------------

\newcommand{\sA}{{\mathscr{A}}}   \newcommand{\sB}{{\mathscr{B}}}
\newcommand{\sC}{{\mathscr{C}}}   \newcommand{\sD}{{\mathscr{D}}}
\newcommand{\sE}{{\mathscr{E}}}   \newcommand{\sF}{{\mathscr{F}}}
\newcommand{\sG}{{\mathscr{G}}}   \newcommand{\sH}{{\mathscr{H}}}
\newcommand{\sI}{{\mathscr{I}}}   \newcommand{\sJ}{{\mathscr{J}}}
\newcommand{\sK}{{\mathscr{K}}}   \newcommand{\sL}{{\mathscr{L}}}
\newcommand{\sM}{{\mathscr{M}}}   \newcommand{\sN}{{\mathscr{N}}}
\newcommand{\sO}{{\mathscr{O}}}   \newcommand{\sP}{{\mathscr{P}}}
\newcommand{\sQ}{{\mathscr{Q}}}   \newcommand{\sR}{{\mathscr{R}}}
\newcommand{\sS}{{\mathscr{S}}}   \newcommand{\sT}{{\mathscr{T}}}
\newcommand{\sU}{{\mathscr{U}}}   \newcommand{\sV}{{\mathscr{V}}}
\newcommand{\sW}{{\mathscr{W}}}   \newcommand{\sX}{{\mathscr{X}}}
\newcommand{\sY}{{\mathscr{Y}}}   \newcommand{\sZ}{{\mathscr{Z}}}


%---------------------------------------
% Math Fraktur Font
%---------------------------------------

%Captital Letters
\newcommand{\mfA}{\mathfrak{A}}	\newcommand{\mfB}{\mathfrak{B}}
\newcommand{\mfC}{\mathfrak{C}}	\newcommand{\mfD}{\mathfrak{D}}
\newcommand{\mfE}{\mathfrak{E}}	\newcommand{\mfF}{\mathfrak{F}}
\newcommand{\mfG}{\mathfrak{G}}	\newcommand{\mfH}{\mathfrak{H}}
\newcommand{\mfI}{\mathfrak{I}}	\newcommand{\mfJ}{\mathfrak{J}}
\newcommand{\mfK}{\mathfrak{K}}	\newcommand{\mfL}{\mathfrak{L}}
\newcommand{\mfM}{\mathfrak{M}}	\newcommand{\mfN}{\mathfrak{N}}
\newcommand{\mfO}{\mathfrak{O}}	\newcommand{\mfP}{\mathfrak{P}}
\newcommand{\mfQ}{\mathfrak{Q}}	\newcommand{\mfR}{\mathfrak{R}}
\newcommand{\mfS}{\mathfrak{S}}	\newcommand{\mfT}{\mathfrak{T}}
\newcommand{\mfU}{\mathfrak{U}}	\newcommand{\mfV}{\mathfrak{V}}
\newcommand{\mfW}{\mathfrak{W}}	\newcommand{\mfX}{\mathfrak{X}}
\newcommand{\mfY}{\mathfrak{Y}}	\newcommand{\mfZ}{\mathfrak{Z}}
%Small Letters
\newcommand{\mfa}{\mathfrak{a}}	\newcommand{\mfb}{\mathfrak{b}}
\newcommand{\mfc}{\mathfrak{c}}	\newcommand{\mfd}{\mathfrak{d}}
\newcommand{\mfe}{\mathfrak{e}}	\newcommand{\mff}{\mathfrak{f}}
\newcommand{\mfg}{\mathfrak{g}}	\newcommand{\mfh}{\mathfrak{h}}
\newcommand{\mfi}{\mathfrak{i}}	\newcommand{\mfj}{\mathfrak{j}}
\newcommand{\mfk}{\mathfrak{k}}	\newcommand{\mfl}{\mathfrak{l}}
\newcommand{\mfm}{\mathfrak{m}}	\newcommand{\mfn}{\mathfrak{n}}
\newcommand{\mfo}{\mathfrak{o}}	\newcommand{\mfp}{\mathfrak{p}}
\newcommand{\mfq}{\mathfrak{q}}	\newcommand{\mfr}{\mathfrak{r}}
\newcommand{\mfs}{\mathfrak{s}}	\newcommand{\mft}{\mathfrak{t}}
\newcommand{\mfu}{\mathfrak{u}}	\newcommand{\mfv}{\mathfrak{v}}
\newcommand{\mfw}{\mathfrak{w}}	\newcommand{\mfx}{\mathfrak{x}}
\newcommand{\mfy}{\mathfrak{y}}	\newcommand{\mfz}{\mathfrak{z}}


\newcommand{\mytitle}{FMA\_3TC13\_TP Poly Exercices 7}
\newcommand{\myauthor}{github/omgeta}
\newcommand{\mydate}{AU 25-26 S1}

\begin{document}
\raggedright
\footnotesize
\begin{center}
{\normalsize{\textbf{\mytitle}}} \\
{\footnotesize{\mydate\hspace{2pt}\textemdash\hspace{2pt}\myauthor}}
\end{center}
\setlist{topsep=-1em, itemsep=-1em, parsep=2em}
\setlist[enumerate]{topsep=0.4em,itemsep=0.8em,parsep=0.35em}
%%%%%%%%%%%%%%%%%%%%%%%%%%%%%%%%%%%%%%%%%%%%%%%%%%%%%%
%                      Begin                         %
%%%%%%%%%%%%%%%%%%%%%%%%%%%%%%%%%%%%%%%%%%%%%%%%%%%%%%
\begin{enumerate}[label=7.\arabic*., start=68]

  \item \begin{minipage}[t]{\linewidth}\raggedright
\textbf{Question.} Justifier le critère de la remarque 7.2 : $f$ est linéaire si et seulement si $f(x+\lambda y)=f(x)+\lambda f(y)$ pour tous $x,y$ et $\lambda\in K$.
\par\smallskip
\textbf{Réponse.} La linéarité implique $f(x+\lambda y)=f(x)+\lambda f(y)$ par additivité puis homogénéité.
    Réciproquement, cette identité donne $f(0)=0$ en prenant $x=y=0$, $\lambda=1$.
    Prendre $\lambda=1$ donne l'additivité, puis $x=0$ donne l'homogénéité.
\end{minipage}

  \item \begin{minipage}[t]{\linewidth}\raggedright
\textbf{Énoncé.} Chaque application est-elle linéaire ? Si oui, donner une matrice dans des bases au choix et dire si c’est un isomorphisme.

\begin{enumerate}[\arabic*.]
        \item \textbf{Question.} $f_1:\RR^2\to\RR^2$, $(x,y)\mapsto(2x+y,x-y)$.
\par\smallskip
\textbf{Réponse.} Linéaire, de matrice $\begin{pmatrix}2&1\\1&-1\end{pmatrix}$ dans les bases canoniques.
        Son déterminant est $-3\neq0$ : c'est un isomorphisme.

  \item \textbf{Question.} $f_2:\RR^3\to\RR^3$, $(x,y,z)\mapsto(xy,x,y)$.
\par\smallskip
\textbf{Réponse.} Non linéaire : $f_2(2,2,0)=(4,2,2)\neq2f_2(1,1,0)$.

  \item \textbf{Question.} $f_3:\RR{}[X]_{\leq3}\to\RR^3$, $P\mapsto(P(-1),P(0),P(1))$.
\par\smallskip
\textbf{Réponse.} Linéaire, de matrice $\begin{pmatrix}1&-1&1&-1\\1&0&0&0\\1&1&1&1\end{pmatrix}$ dans les bases canoniques.
        Son rang vaut $3$, son noyau est $\RR(X^3-X)$.
        Elle est surjective, mais n'est pas un isomorphisme.

  \item \textbf{Question.} $f_4:C^0(\RR,\RR)\to\RR$, $h\mapsto\int_0^1h(x)^2dx$.
\par\smallskip
\textbf{Réponse.} Non linéaire : pour la fonction constante $1$, $f_4(2)=4\neq2f_4(1)$.

\end{enumerate}
\end{minipage}

  \item \begin{minipage}[t]{\linewidth}\raggedright
\textbf{Question.} Soit $A\in\RR^{2\times2}$. Montrer que $\psi:M\mapsto AM-MA$ est linéaire et donner sa matrice dans la base $(E_{11},E_{12},E_{21},E_{22})$.
\par\smallskip
\textbf{Réponse.} $\psi(\lambda M+N)=\lambda(AM-MA)+(AN-NA)$, donc $\psi$ est linéaire.
    Pour $A=\begin{pmatrix}a&b\\c&d\end{pmatrix}$, sa matrice dans $(E_{11},E_{12},E_{21},E_{22})$ est
    \[
      \begin{pmatrix}0&-c&b&0\\-b&a-d&0&b\\c&0&d-a&-c\\0&c&-b&0\end{pmatrix}
    \]
\end{minipage}

  \item \begin{minipage}[t]{\linewidth}\raggedright
\textbf{Question.} Soient $\dim E=n$, $\varphi^n=0$, $\varphi^{n-1}\neq0$ et $x$ tel que $\varphi^{n-1}(x)\neq0$. Donner un exemple et montrer que $(x,\varphi(x),\ldots,\varphi^{n-1}(x))$ est une base de $E$.
\par\smallskip
\textbf{Réponse.} Exemple : $\varphi(e_i)=e_{i+1}$ pour $i<n$, $\varphi(e_n)=0$, et $x=e_1$.
    Pour montrer la liberté en général, supposons $\sum_{k=0}^{n-1}\lambda_k\varphi^k(x)=0$.
    Si $r$ est le plus petit indice de coefficient non nul, appliquer $\varphi^{n-1-r}$ donne
    $\lambda_r\varphi^{n-1}(x)=0$, contradiction.
    $\therefore$ la famille est libre de cardinal $n$, donc c'est une base.
\end{minipage}

  \item \begin{minipage}[t]{\linewidth}\raggedright
\textbf{Énoncé.} Soient $E_1,E_2$ deux sous-espaces de dimension finie de $E$ et $f:E_1\times E_2\to E$, $(x_1,x_2)\mapsto x_1+x_2$.

\begin{enumerate}[\arabic*.]
        \item \textbf{Question.} Montrer que $f$ est linéaire.
\par\smallskip
\textbf{Réponse.} $f(\lambda(x_1,x_2)+(y_1,y_2))=\lambda(x_1+x_2)+y_1+y_2$, donc $f$ est linéaire.

  \item \textbf{Question.} Déterminer son noyau et son image.
\par\smallskip
\textbf{Réponse.} $\ker f=\{(x,-x):x\in E_1\cap E_2\}$ et $\operatorname{im}f=E_1+E_2$.

  \item \textbf{Question.} Que donne le théorème du rang ?
\par\smallskip
\textbf{Réponse.} Le théorème du rang donne $\dim E_1+\dim E_2=\dim(E_1\cap E_2)+\dim(E_1+E_2)$.

\end{enumerate}
\end{minipage}

  \item \begin{minipage}[t]{\linewidth}\raggedright
\textbf{Énoncé.} Soit $E=F_1\oplus\cdots\oplus F_n$ de dimension finie. On note $L(E,E)$ l’espace des endomorphismes.

\begin{enumerate}[\arabic*.]
        \item \textbf{Question.} Pour $\mathcal F_i=\{u:\operatorname{im}u\subseteq F_i\}$, montrer que $L(E,E)=\mathcal F_1\oplus\cdots\oplus\mathcal F_n$.
\par\smallskip
\textbf{Réponse.} Soient $p_i$ les projections associées à $E=\bigoplus_iF_i$.

Tout $u$ s'écrit $u=\sum_i p_i u$, avec $\operatorname{im}(p_i u)\subseteq F_i$.
        L'unicité vient de celle de la décomposition de chaque $u(x)$ dans les $F_i$.
        $\therefore L(E,E)=\bigoplus_i\mathcal F_i$.

  \item \textbf{Question.} Pour $\mathcal G_i=\{u:F_i\subseteq\ker u\}$, montrer que $L(E,E)=\mathcal G_1\oplus\cdots\oplus\mathcal G_n$ (énoncé imprimé).
\par\smallskip
\textbf{Réponse.} L'énoncé imprimé est faux en général pour $n\geq3$.
        Dans $E=K^3$, $F_i=Ke_i$, une application non nulle sur $e_3$ et nulle sur $e_1,e_2$ appartient à $\mathcal G_1\cap\mathcal G_2$.

        La version correcte est $\widetilde{\mathcal G}_i=\{u:\bigoplus_{j\neq i}F_j\subseteq\ker u\}$.
        Alors $u=\sum_i up_i$ est l'unique décomposition, donc $L(E,E)=\bigoplus_i\widetilde{\mathcal G}_i$.
        Pour $n=2$, la formule imprimée est vraie, les deux indices étant échangés. Pour $n=1$ et $E\neq0$, elle est fausse car $\mathcal G_1=\{0\}$.

\end{enumerate}
\end{minipage}

  \item \begin{minipage}[t]{\linewidth}\raggedright
\textbf{Question.} Soit $f\in L(E,E)$ avec $\dim E=n$. Montrer que $\ker f=\operatorname{im}f$ équivaut à $f^2=0$ et $n=2\operatorname{rg}f$.
\par\smallskip
\textbf{Réponse.} Si $\ker f=\operatorname{im}f$, alors $f^2=0$ et le théorème du rang donne $n=2\operatorname{rg}f$.
    Réciproquement, $f^2=0$ donne $\operatorname{im}f\subseteq\ker f$.
    Si $n=2\operatorname{rg}f$, les deux espaces ont la même dimension et sont donc égaux.
\end{minipage}

  \item \begin{minipage}[t]{\linewidth}\raggedright
\textbf{Énoncé.} Soient $f:E\to F$ et $g:F\to G$ linéaires.

\begin{enumerate}[\arabic*.]
        \item \begin{enumerate}[a.]
            \item \textbf{Question.} Exprimer $g\circ f=0$ en termes de noyau et d’image.
\par\smallskip
\textbf{Réponse.} $g\circ f=0\iff\operatorname{im}f\subseteq\ker g$.

  \item \textbf{Question.} En dimension finie, quelle relation en déduire entre $\operatorname{rg}f$ et $\operatorname{rg}g$ ?
\par\smallskip
\textbf{Réponse.} Donc $\operatorname{rg}f\leq\dim\ker g=\dim F-\operatorname{rg}g$, soit $\operatorname{rg}f+\operatorname{rg}g\leq\dim F$.

\end{enumerate}

  \item \textbf{Question.} Montrer que $f(\ker(g\circ f))=\ker g\cap\operatorname{im}f$.
\par\smallskip
\textbf{Réponse.} $y\in f(\ker(gf))\iff\exists x,\ y=f(x),\ g(f(x))=0
        \iff y\in\operatorname{im}f\cap\ker g$.

\end{enumerate}
\end{minipage}

  \item \begin{minipage}[t]{\linewidth}\raggedright
\textbf{Énoncé.} Soient $u,v:E\to F$ linéaires, avec $E,F$ de dimension finie.

\begin{enumerate}[\arabic*.]
        \item \textbf{Question.} Montrer que $\operatorname{rg}(u+v)\leq\operatorname{rg}u+\operatorname{rg}v$.
\par\smallskip
\textbf{Réponse.} $\operatorname{im}(u+v)\subseteq\operatorname{im}u+\operatorname{im}v$, donc $\operatorname{rg}(u+v)\leq\operatorname{rg}u+\operatorname{rg}v$.

  \item \textbf{Question.} En déduire $|\operatorname{rg}u-\operatorname{rg}v|\leq\operatorname{rg}(u+v)$.
\par\smallskip
\textbf{Réponse.} Appliquer la question 1 à $u=(u+v)-v$ donne $\operatorname{rg}u-\operatorname{rg}v\leq\operatorname{rg}(u+v)$.
        Échanger $u,v$ donne l'autre inégalité.

\end{enumerate}
\end{minipage}

  \item \begin{minipage}[t]{\linewidth}\raggedright
\textbf{Énoncé.} Dans $V=\RR^3$, soient $u=(4,-2,1)$, $w_1=(6,-5,1)$, $w_2=(-5,8,0)$, $U=\RR u$ et $W=\operatorname{Vect}(w_1,w_2)$.

\begin{enumerate}[\arabic*.]
        \item \textbf{Question.} Vérifier que $V=U\oplus W$.
\par\smallskip
\textbf{Réponse.} Pour $P=(u\ w_1\ w_2)$, $\det P=1\neq0$.
        $\therefore V=U\oplus W$.

  \item \textbf{Question.} Donner dans la base canonique la matrice de la projection sur $W$ parallèlement à $U$.
\par\smallskip
\textbf{Réponse.} $p_{W,U}=P\operatorname{diag}(0,1,1)P^{-1}$, donc
        \[
          [p_{W,U}]=\begin{pmatrix}33&20&-92\\-16&-9&46\\8&5&-22\end{pmatrix}
        \]

\end{enumerate}
\end{minipage}

  \item \begin{minipage}[t]{\linewidth}\raggedright
\textbf{Question.} Dans $\RR^3$, soient $u_1=(1,-1,0)$, $u_2=(-1,0,-1)$, $w=(4,3,6)$, $U=\operatorname{Vect}(u_1,u_2)$ et $W=\RR w$. Donner la matrice canonique de la symétrie par rapport à $W$ selon $U$.
\par\smallskip
\textbf{Réponse.} Pour $P=(u_1\ u_2\ w)$, la matrice cherchée est
    \[
      S=P\operatorname{diag}(-1,-1,1)P^{-1}
       =\begin{pmatrix}7&8&-8\\6&5&-6\\12&12&-13\end{pmatrix}
    \]
    On vérifie $Su_i=-u_i$ et $Sw=w$.
\end{minipage}

  \item \begin{minipage}[t]{\linewidth}\raggedright
\textbf{Question.} Soit $E$ un espace vectoriel. Montrer que $f:E\to E$ est une homothétie si et seulement si $(x,f(x))$ est liée pour tout $x\in E$ (la linéarité de $f$ est implicite).
\par\smallskip
\textbf{Réponse.} Il faut supposer $f$ linéaire : sinon $f(x)=x^2$ sur $\RR$ donne un contre-exemple.
    Sous cette hypothèse, si $f=\lambda\operatorname{id}$, $(x,f(x))$ est liée pour tout $x$.
    Réciproquement, pour $x\neq0$, écrivons $f(x)=\lambda_xx$.
    Si $x,y$ sont indépendants, comparer $f(x+y)$ à $f(x)+f(y)$ donne
    $\lambda_x=\lambda_y=\lambda_{x+y}$.
    Si $x,y$ sont colinéaires non nuls, la linéarité donne aussi $\lambda_x=\lambda_y$.
    $\therefore f$ est une homothétie.
\end{minipage}

  \item \begin{minipage}[t]{\linewidth}\raggedright
\textbf{Énoncé.} Soient $\dim E=3$ sur $\RR$ ou $\CC$, et $f\in L(E,E)$.

\begin{enumerate}[\arabic*.]
        \item \textbf{Question.} Si $f\neq0$ et $f^2=0$, montrer qu’il existe une base où $[f]=\begin{pmatrix}0&0&0\\1&0&0\\0&0&0\end{pmatrix}$.
\par\smallskip
\textbf{Réponse.} $\operatorname{im}f\subseteq\ker f$ donne $\operatorname{rg}f\leq1$.
        Comme $f\neq0$, le rang est $1$ et $\dim\ker f=2$.
        Choisir $x$ avec $f(x)\neq0$, puis $z\in\ker f\setminus Kf(x)$.
        Dans $(x,f(x),z)$, la matrice est $\begin{pmatrix}0&0&0\\1&0&0\\0&0&0\end{pmatrix}$.

  \item \textbf{Question.} Si $f^2\neq0$ et $f^3=0$, montrer qu’il existe une base où $[f]=\begin{pmatrix}0&0&0\\1&0&0\\0&1&0\end{pmatrix}$.
\par\smallskip
\textbf{Réponse.} Choisir $x$ avec $f^2(x)\neq0$.
        D'après 7.71, $(x,f(x),f^2(x))$ est une base.
        La matrice est $\begin{pmatrix}0&0&0\\1&0&0\\0&1&0\end{pmatrix}$.

\end{enumerate}
\end{minipage}

  \item \begin{minipage}[t]{\linewidth}\raggedright
\textbf{Question.} Soient $\dim E=3$ sur $\RR$ ou $\CC$, $f^2\neq0$, $f^3=0$, et $g$ commutant avec $f$. Montrer que $g$ est une combinaison linéaire de $I,f,f^2$.
\par\smallskip
\textbf{Réponse.} Choisir $x$ tel que $(x,f(x),f^2(x))$ soit une base.
    Écrire $g(x)=ax+bf(x)+cf^2(x)$.
    Comme $gf=fg$, on a aussi $g(f(x))=(aI+bf+cf^2)f(x)$ et de même sur $f^2(x)$.
    Les deux applications coïncident sur une base.
    $\therefore g=aI+bf+cf^2$.
\end{minipage}

  \item \begin{minipage}[t]{\linewidth}\raggedright
\textbf{Énoncé.} On travaille sur $K=\RR$ ou $\CC$.

\begin{enumerate}[\arabic*.]
        \item \textbf{Question.} Si $fg-gf=I$, montrer que $f^ng-gf^n=nf^{n-1}$ pour tout $n\geq1$.
\par\smallskip
\textbf{Réponse.} Pour $n=1$, c'est l'hypothèse. Si $f^ng-gf^n=nf^{n-1}$,
        \[
          f^{n+1}g-gf^{n+1}=f(f^ng-gf^n)+(fg-gf)f^n=(n+1)f^n
        \]

  \item \textbf{Question.} Montrer qu’une telle paire est impossible si $E$ est de dimension finie non nulle.
\par\smallskip
\textbf{Réponse.} En dimension finie non nulle, $\operatorname{tr}(fg-gf)=0$, alors que $\operatorname{tr}I=\dim E\neq0$ sur $\RR$ ou $\CC$.
        Contradiction.

  \item \textbf{Question.} Dans $K[X]$, examiner $f:P\mapsto P^{\prime}$ et $g:P\mapsto XP$.
\par\smallskip
\textbf{Réponse.} Sur $K[X]$, $f(g(P))-g(f(P))=(XP)'-XP'=P$.
        C'est bien un exemple en dimension infinie.

\end{enumerate}
\end{minipage}

  \item \begin{minipage}[t]{\linewidth}\raggedright
\textbf{Question.} Calculer le déterminant de la conjugaison complexe $\chi:z\mapsto\overline z$, considérée comme endomorphisme du $\RR$-espace vectoriel $\CC$.
\par\smallskip
\textbf{Réponse.} Dans la base réelle $(1,i)$ de $\CC$, $[\chi]=\operatorname{diag}(1,-1)$.
    $\therefore\det\chi=-1$.
\end{minipage}

  \item \begin{minipage}[t]{\linewidth}\raggedright
\textbf{Énoncé.} Soit $E$ réel de dimension $n\geq2$ et $u\in L(E,E)$ tel que $u^2=-I$.

\begin{enumerate}[\arabic*.]
        \item \textbf{Question.} Montrer que $n$ est pair.
\par\smallskip
\textbf{Réponse.} $(\det u)^2=\det(-I)=(-1)^n$.
        Le membre de gauche est strictement positif, donc $n$ est pair.

  \item \textbf{Question.} Montrer que $u$ ne laisse stable aucun hyperplan de $E$.
\par\smallskip
\textbf{Réponse.} Si un hyperplan $H$ était stable, $u$ induirait sur $E/H$ un endomorphisme $\lambda I$ vérifiant $\lambda^2=-1$.
        Impossible sur $\RR$.

\end{enumerate}
\end{minipage}

  \item \begin{minipage}[t]{\linewidth}\raggedright
\textbf{Énoncé.} Soit $a$ l’endomorphisme de $\RR^4$ de matrice canonique $\displaystyle A=\begin{pmatrix}1&1&1&1\\1&0&0&1\\1&0&0&1\\1&1&1&1\end{pmatrix}$.

\begin{enumerate}[\arabic*.]
        \item \textbf{Question.} Donner une base de $\operatorname{im}a$ et de $\ker a$.
\par\smallskip
\textbf{Réponse.} Posons $r=(1,1,1,1)^\top$, $s=(1,0,0,1)^\top$,
    $k_1=(1,0,0,-1)^\top$, $k_2=(0,1,-1,0)^\top$.

$\operatorname{im}a=\operatorname{Vect}(r,s)$ et $\ker a=\operatorname{Vect}(k_1,k_2)$.

  \item \textbf{Question.} Montrer que $\operatorname{im}a$ est stable par $a$.
\par\smallskip
\textbf{Réponse.} L'image de tout endomorphisme est stable : $a(a(x))\in\operatorname{im}a$.

  \item \textbf{Question.} Donner la matrice de $\varphi=a|_{\operatorname{im}a}$ dans une base de l’image.
\par\smallskip
\textbf{Réponse.} Dans $(r,s)$, $a(r)=2r+2s$ et $a(s)=2r$, donc $[\varphi]=\begin{pmatrix}2&2\\2&0\end{pmatrix}$.

  \item \textbf{Question.} Montrer que $A$ est semblable à $\operatorname{diag}\left(0_2,\begin{pmatrix}2&2\\2&0\end{pmatrix}\right)$.
\par\smallskip
\textbf{Réponse.} $(k_1,k_2,r,s)$ est une base (l'intersection noyau-image est nulle).
        La matrice dans cette base est exactement $B$.

\end{enumerate}
\end{minipage}

  \item \begin{minipage}[t]{\linewidth}\raggedright
\textbf{Question.} Soient $p$ un projecteur et $u$ un endomorphisme de $E$. Montrer que $pu=up$ si et seulement si $\ker p$ et $\operatorname{im}p$ sont stables par $u$.
\par\smallskip
\textbf{Réponse.} Si $pu=up$, $p(x)=0\Rightarrow p(u(x))=0$, et $x=p(y)\Rightarrow u(x)=p(u(y))$.
    Ainsi noyau et image sont stables.
    Réciproquement, écrire $x=y+z$ avec $y\in\operatorname{im}p$, $z\in\ker p$.
    La stabilité donne $pu(x)=u(y)=up(x)$.
\end{minipage}

  \item \begin{minipage}[t]{\linewidth}\raggedright
\textbf{Question.} Soient $\lambda\neq0$, $\mu\in K$ et $P_{n+1}=\lambda P_n+\mu P_n^{\prime}$ dans $K[X]_{\leq2}$. Pour $n_0\in\NN$ et $Q\in K[X]_{\leq2}$, existe-t-il $P_0$ tel que $P_{n_0}=Q$ ?
\par\smallskip
\textbf{Réponse.} Sur $K[X]_{\leq2}$, $T=\lambda I+\mu D$ est triangulaire de diagonale $\lambda\neq0$, donc inversible.
    Comme $D^3=0$,
    $T^{-1}=\lambda^{-1}(I-(\mu/\lambda)D+(\mu/\lambda)^2D^2)$.
    $\therefore$ il existe un unique $P_0=T^{-n_0}Q$ (pour $n_0\in\NN$).
\end{minipage}

  \item \begin{minipage}[t]{\linewidth}\raggedright
\textbf{Question.} Soit $f$ un endomorphisme d’un espace $V$ de dimension finie. Montrer l’équivalence de $V=\operatorname{im}f\oplus\ker f$, $\operatorname{im}f=\operatorname{im}f^2$, et $\ker f=\ker f^2$.
\par\smallskip
\textbf{Réponse.} $\operatorname{im}f^2\subseteq\operatorname{im}f$ et $\ker f\subseteq\ker f^2$.
    En dimension finie, leurs égalités sont équivalentes par le théorème du rang.
    De plus,
    \[
      \ker f^2=\ker f\iff\operatorname{im}f\cap\ker f=\{0\}
    \]
    En effet, $f(x)\in\ker f$ implique $x\in\ker f^2$ ; la réciproque suit de la même observation.
    Les dimensions de l'image et du noyau ont pour somme $\dim V$, donc les trois assertions sont équivalentes.
\end{minipage}

  \item \begin{minipage}[t]{\linewidth}\raggedright
\textbf{Question.} Soient $A,B\in\CC^{n\times n}$ avec $A^2B=A$ et $\operatorname{rg}A=\operatorname{rg}B$. Montrer que $B^2A=B$.
\par\smallskip
\textbf{Réponse.} $A^2B=A$ donne $\ker B\subseteq\ker A$, donc $\ker B=\ker A$ par égalité des rangs.
    Elle donne aussi $\operatorname{im}A^2=\operatorname{im}A$, donc $\CC^n=\operatorname{im}A\oplus\ker A$.
    Dans une base adaptée,
    \[
      A=\begin{pmatrix}C&0\\0&0\end{pmatrix},\quad
      B=\begin{pmatrix}C^{-1}&0\\D&0\end{pmatrix},\quad C\text{ inversible}
    \]
    Un produit direct donne $B^2A=B$.
\end{minipage}

  \item \begin{minipage}[t]{\linewidth}\raggedright
\textbf{Question.} Soit $f\in L(V,V)$ en dimension finie avec $V=\operatorname{im}f\oplus\ker f$. Posons $m=\dim\operatorname{im}f$, $m^{\prime}=\dim\ker f$. Montrer qu’une base donne $[f]=\operatorname{diag}(0_{m^{\prime}},A)$ avec $A\in\operatorname{GL}_m(K)$.
\par\smallskip
\textbf{Réponse.} Prendre une base du noyau, puis une base de l'image.
    La restriction de $f$ à son image est injective, car $\operatorname{im}f\cap\ker f=0$.
    Elle est donc inversible en dimension finie.
    Dans la base réunie, $[f]=\operatorname{diag}(0_{m'},A)$ avec $A$ inversible.
\end{minipage}

  \item \begin{minipage}[t]{\linewidth}\raggedright
\textbf{Énoncé.} Sur $E=K[X]_{\leq n}$, on définit $f(P)=P(X+1)+P(X-1)-2P(X)$.

\begin{enumerate}[\arabic*.]
        \item \textbf{Question.} Montrer que $f$ est un endomorphisme.
\par\smallskip
\textbf{Réponse.} Sur un corps de caractéristique $0$ (ou strictement supérieure à $n$) :

Translations et soustraction étant linéaires, $f$ est un endomorphisme.

  \item \textbf{Question.} Calculer $f(X^p)$ pour $0\leq p\leq n$ et son degré ; en déduire noyau, image et rang.
\par\smallskip
\textbf{Réponse.} $\displaystyle f(X^p)=2\sum_{j=1}^{\lfloor p/2\rfloor}\binom p{2j}X^{p-2j}$.
        Pour $p\geq2$, son degré est $p-2$, de coefficient dominant $p(p-1)$.
        Donc $\ker f=K[X]_{\leq1}$, $\operatorname{im}f=K[X]_{\leq n-2}$, $\operatorname{rg}f=n-1$ pour $n\geq2$.
        Pour $n\leq1$, $f=0$.

  \item \textbf{Question.} Pour $Q\in\operatorname{im}f$, montrer qu’il existe un unique $P\in E$ tel que $f(P)=Q$ et $P(0)=P^{\prime}(0)=0$.
\par\smallskip
\textbf{Réponse.} Si $P_1$ est un antécédent de $Q$, poser $P=P_1-P_1(0)-P_1'(0)X$.
        Il vérifie les deux conditions et $f(P)=Q$.
        La différence de deux solutions appartient au noyau et a ses deux premiers coefficients nuls : elle est nulle.

En caractéristique $p\leq n$, la question 3 est fausse : $X^p$ est dans le noyau et vérifie $P(0)=P'(0)=0$.
    Pour préciser le noyau et le rang, poser $H=X^p-X$ et $d_q=\min(p,n-pq+1)$, $0\leq q\leq\lfloor n/p\rfloor$.
    Les blocs $H^q\operatorname{Vect}(1,X,\ldots,X^{d_q-1})$ sont stables.
    Si $p>2$, le noyau a pour base les $H^q$, ainsi que les $XH^q$ lorsque $d_q\geq2$ ; l'image a pour base les $X^rH^q$ avec $0\leq r\leq d_q-3$.
    Donc $\operatorname{rg}f=\sum_q\max(d_q-2,0)$.
    Si $p=2$, $f=0$ car $X+1=X-1$.

\end{enumerate}
\end{minipage}

  \item \begin{minipage}[t]{\linewidth}\raggedright
\textbf{Question.} Pour $u:E\to F$, $v:F\to G$ linéaires, montrer que $\ker(v\circ u)=u^{-1}(\ker v)$.
\par\smallskip
\textbf{Réponse.} $x\in\ker(vu)\iff v(u(x))=0\iff u(x)\in\ker v\iff x\in u^{-1}(\ker v)$.
\end{minipage}

  \item \begin{minipage}[t]{\linewidth}\raggedright
\textbf{Question.} Soient $u,v\in L(E,E)$ avec $u^2=u$ et $vu=0$. Montrer que $\operatorname{im}(u+v)=\operatorname{im}u+\operatorname{im}v$.
\par\smallskip
\textbf{Réponse.} $(u+v)u=u^2+vu=u$, donc $\operatorname{im}u\subseteq\operatorname{im}(u+v)$.
    Puis $v(x)=(u+v)(x)-u(x)$ donne $\operatorname{im}v\subseteq\operatorname{im}(u+v)$.
    L'autre inclusion est immédiate.
    $\therefore\operatorname{im}(u+v)=\operatorname{im}u+\operatorname{im}v$.
\end{minipage}

  \item \begin{minipage}[t]{\linewidth}\raggedright
\textbf{Question.} Soit $f$ un endomorphisme d’un espace de dimension finie. Montrer qu’il existe deux isomorphismes $u,v$ tels que $f=u-v$.
\par\smallskip
\textbf{Réponse.} Sur un corps infini (en particulier $\RR$ ou $\CC$), choisir $t\neq0$ tel que $\det(f+tI)\neq0$ : il n'y a qu'un nombre fini de valeurs exclues.
    Alors $f=(f+tI)-tI$ est une différence d'isomorphismes.
    Sans hypothèse sur le corps, l'énoncé est faux : sur $\ZZ/2\ZZ$ en dimension $1$, l'identité n'est pas une différence de deux automorphismes.
\end{minipage}

  \item \begin{minipage}[t]{\linewidth}\raggedright
\textbf{Question.} À quelle condition $K[X]_{\leq m}$ est-il un hyperplan de $K[X]_{\leq n}$ ?
\par\smallskip
\textbf{Réponse.} $K[X]_{\leq m}$ est inclus dans $K[X]_{\leq n}$ si et seulement si $m\leq n$.
    La codimension est alors $(n+1)-(m+1)=n-m$.
    $\therefore$ c'est un hyperplan si et seulement si $m=n-1$.
\end{minipage}

  \item \begin{minipage}[t]{\linewidth}\raggedright
\textbf{Question.} Soient $f,g$ deux formes linéaires de même noyau. Montrer qu’il existe $\alpha\neq0$ tel que $g=\alpha f$.
\par\smallskip
\textbf{Réponse.} Si $f=0$, alors $g=0$ et $\alpha=1$ convient.
    Sinon choisir $x_0$ tel que $f(x_0)=1$.
    Pour tout $x$, $x-f(x)x_0\in\ker f=\ker g$, donc $g(x)=g(x_0)f(x)$.
    Comme $x_0\notin\ker g$, $\alpha=g(x_0)\neq0$.
\end{minipage}

  \item \begin{minipage}[t]{\linewidth}\raggedright
\textbf{Énoncé.} Soient $n\geq2$ et $H$ un hyperplan de $K^{n\times n}$.

\begin{enumerate}[\arabic*.]
        \item \textbf{Question.} Montrer qu’il existe $A\neq0$ telle que $H=\{M:\operatorname{tr}(AM)=0\}$.
\par\smallskip
\textbf{Réponse.} $H=\ker\ell$ pour une forme linéaire non nulle.
        Poser $a_{i,j}=\ell(E_{j,i})$ donne $\ell(M)=\operatorname{tr}(AM)$ et $A\neq0$.

  \item \textbf{Question.} En déduire que $H$ contient une matrice inversible.
\par\smallskip
\textbf{Réponse.} Il existe $U,V$ inversibles tels que $UAV=\operatorname{diag}(I_r,0)$.
        Prendre la matrice $P$ de la permutation cyclique de $\{1,\ldots,n\}$, sans point fixe.
        Elle est inversible et sa diagonale est nulle.
        Alors $M=VPU$ est inversible et
        $\operatorname{tr}(AM)=\operatorname{tr}(UAVP)=0$.
        $\therefore M\in H$.

\end{enumerate}
\end{minipage}

  \item \begin{minipage}[t]{\linewidth}\raggedright
\textbf{Énoncé.} Pour $n\geq2$, on définit $\varphi:K[X]_{\leq n}\to K[X]_{\leq n}$ par $\varphi(P)=X^nP(1/X)$.

\begin{enumerate}[\arabic*.]
        \item \textbf{Question.} Vérifier la linéarité de $\varphi$.
\par\smallskip
\textbf{Réponse.} $\varphi(\sum_{k=0}^na_kX^k)=\sum_{k=0}^na_kX^{n-k}$ : les coefficients sont renversés linéairement.

  \item \textbf{Question.} Calculer $\varphi\circ\varphi$.
\par\smallskip
\textbf{Réponse.} Renverser deux fois donne $\varphi^2=I$.

  \item \textbf{Question.} Reconnaître géométriquement $\varphi$.
\par\smallskip
\textbf{Réponse.} Si $\operatorname{car}K\neq2$, c'est la symétrie par rapport aux polynômes à coefficients symétriques ($a_k=a_{n-k}$), selon ceux à coefficients antisymétriques ($a_k=-a_{n-k}$).
        En caractéristique $2$, c'est toujours une involution, mais la décomposition en ces deux espaces ne s'applique pas.

\end{enumerate}
\end{minipage}

  \item \begin{minipage}[t]{\linewidth}\raggedright
\textbf{Question.} Soient $A,B$ réelles avec $A^2=B^2=0_{2n}$ et $\operatorname{rg}A,\operatorname{rg}B\geq n$. Montrer qu’elles sont semblables. Le poly affiche $n\times n$, mais les identités imposent ici la taille $2n\times2n$.
\par\smallskip
\textbf{Réponse.} La taille cohérente est $2n\times2n$, puisque $A^2=B^2=0_{2n}$.
    $\operatorname{im}A\subseteq\ker A$ donne $\operatorname{rg}A\leq n$, donc $\operatorname{rg}A=n$.
    Choisir $u_1,\ldots,u_n$ tels que les $Au_i$ forment une base de l'image.
    Appliquer $A$ à une relation montre que $(u_1,\ldots,u_n,Au_1,\ldots,Au_n)$ est libre, donc une base.
    Dans cette base, $[A]=\begin{pmatrix}0&0\\I_n&0\end{pmatrix}$.
    Le même raisonnement vaut pour $B$.
    $\therefore A$ et $B$ sont semblables.
\end{minipage}

\end{enumerate}
%%%%%%%%%%%%%%%%%%%%%%%%%%%%%%%%%%%%%%%%%%%%%%%%%%%%%%
%                       End                          %
%%%%%%%%%%%%%%%%%%%%%%%%%%%%%%%%%%%%%%%%%%%%%%%%%%%%%%

\end{document}
